

% ========================================
\subsection{Themenschwerpunkt: Audio I/O und 3D-Raumdaten}
% ========================================

Dieser Abschnitt beschreibt die Implementierung zweier Komponenten des Multithreaded Acoustic Ray Tracers:

\begin{enumerate}
    \item \textbf{Audio I/O:} Einlesen trockener WAV-Dateien und Export des verarbeiteten, halligen Audio-Signals mit Clipping-Pr\"avention
    \item \textbf{Raumdaten:} Laden von 3D-Modellen aus \texttt{.obj}-Dateien und Konvertierung in das interne JSON-Format
\end{enumerate}

Die Implementierung erfolgt in Python und integriert sich nahtlos in die bestehende Pipeline, die Rust-basierte Physik-Engine, DSP-Convolution und Multithreading verbindet.

\subsection{Architektur-\"Ubersicht}

Das System folgt einer erweiterten Pipeline mit automatischer Raumgeometrie-Integration:

\begin{enumerate}
    \item[\textbf{[0/4]}] OBJ-Import: \texttt{room.obj} $\to$ \texttt{room\_geometry.json} mit Material-Absorption
    \item[\textbf{[1/4]}] Rust Physics Engine: L\"adt Raumgeometrie, simuliert Schallstrahlen, erzeugt IR
    \item[\textbf{[2/4]}] Audio-Einlesung: L\"adt trockene WAV-Datei
    \item[\textbf{[3/4]}] DSP-Convolution: Wendet IR mit frequenzabh\"angiger Material-D\"ampfung an
    \item[\textbf{[4/4]}] Audio-Export: Normalisiert und speichert Ergebnis
\end{enumerate}

Das Wandmaterial wird in \texttt{input\_config.json} konfiguriert und durchgehend durch alle Stufen propagiert: von der Geometrie-Absorption im Rust-Tracer bis zur oktavbandbasierten Frequenzformung in der Convolution.

% ========================================
\subsection{Audio I/O Handling}
% ========================================

\subsubsection{Anforderungen (Akzeptanzkriterien)}

\begin{enumerate}
    \item Sauberes Einlesen der \texttt{unprocessed.wav} ohne Qualit\"atsverlust
    \item Export der finalen WAV-Datei ohne Artefakte
    \item Normalisierung, damit geschichtetes Audio beim Export nicht clippt
\end{enumerate}

\subsection{Theoretischer Hintergrund}

\subsubsection{WAV-Format und Audio-Digitalisierung}

Das WAV-Format speichert Audiodaten als zeitdiskrete Samples mit konstanter Abtastfrequenz $f_s$:

$$x[n] = \int_{(n-0.5)/f_s}^{(n+0.5)/f_s} x_{\text{analog}}(t) \, dt$$

Jedes Sample liegt im Bereich $[-1.0, 1.0]$ (bei normalisierten 16/32-Bit PCM).

\subsubsection{WAV-Format: Technische Spezifikation}

Das WAV-Format (Waveform Audio File Format) speichert unkomprimierte PCM-Audiodaten
in einem RIFF-Container. Die f\"ur die Pipeline relevanten Parameter:

\begin{table}[H]
    \centering
    \begin{tabular}{|l|l|p{5cm}|}
        \hline
        \textbf{Parameter} & \textbf{Typische Werte} & \textbf{Bedeutung} \\
        \hline
        Abtastrate $f_s$ & 44.100\,Hz, 48.000\,Hz & Samples pro Sekunde \\
        \hline
        Bit-Tiefe & 16-bit, 24-bit, 32-bit float & Amplitude-Aufl\"osung \\
        \hline
        Kan\"ale & 1 (Mono), 2 (Stereo) & Anzahl unabh\"angiger Signale \\
        \hline
        Kodierung & PCM (Integer) oder IEEE Float & Speicherformat der Samples \\
        \hline
    \end{tabular}
    \caption{WAV-Formatparameter und ihre Bedeutung f\"ur die Pipeline}
\end{table}

Die \texttt{soundfile}-Bibliothek normalisiert beim Einlesen automatisch:
16-bit-Integer-Samples im Bereich $[-32768,\,32767]$ werden in \texttt{float64}-Werte
im Bereich $[-1.0,\,1.0]$ umgewandelt. Diese Skalierung ist verlustfrei -- die
Quantisierungsschritte des PCM-Formats bleiben erhalten. 32-bit-Float-WAV
wird direkt \"ubernommen. Die Pipeline erwartet stets Float-Daten; \texttt{write\_wav}
schreibt standardm\"a\ss{}ig 32-bit-Float zur\"uck, was f\"ur Hall-Anwendungen
bei $f_s = 44.100$\,Hz ausreicht.

\textbf{Abtastraten-Konsistenz:} Die Rust-Engine ist in \texttt{input\_config.json}
mit einer festen Abtastrate konfiguriert ($44.100$\,Hz).
Stimmt die Abtastrate der Eingangs-WAV nicht \"uberein, gibt die Pipeline eine Warnung aus
ohne automatisches Resampling, da eine Rateninkonsistenz alle Verzögerungen
der Impulsantwort um den Ratenquotienten staucht oder dehnt und den Hall
hörbar verstimmt.

\subsubsection{Clipping und Normalisierung}

Wenn w\"ahrend der Convolution (Faltung) die Amplitude $|y[n]|$ \"uber 1.0 steigt, tritt \textbf{Clipping} auf:

$$y_{\text{clipped}}[n] = \begin{cases}
    1.0 & \text{wenn } y[n] > 1.0 \\
    y[n] & \text{wenn } |y[n]| \leq 1.0 \\
    -1.0 & \text{wenn } y[n] < -1.0
\end{cases}$$

Um dies zu vermeiden, normalisieren wir:

$$y_{\text{norm}}[n] = \frac{y[n]}{A_{\max}} \quad \text{wobei} \quad A_{\max} = \max_n |y[n]|$$

Diese Normalisierung ist bewusst \textbf{konditional}: Liegt das Maximum bereits
unter 1.0, wird nicht skaliert da leise Signale nicht k\"unstlich lauter gemacht werden.
Die Pegelrelation zwischen verschiedenen Material-L\"aufen bleibt damit
erhalten, was akustische Vergleiche erm\"oglicht.

\subsection{Implementierung: \texttt{wav\_handler.py}}

\subsubsection{WAV-Einlesung}

\begin{lstlisting}
import numpy as np
import soundfile as sf

def read_wav(wav_path: str) -> tuple:
    """
    Reads a WAV file and returns the audio data as a NumPy array
    along with the corresponding sample rate.
    """
    audio_data, sample_rate = sf.read(wav_path)
    return audio_data, sample_rate
\end{lstlisting}

\textbf{Funktionsweise:}

\begin{itemize}
    \item Die \texttt{soundfile}-Bibliothek liest WAV-Dateien robust
    \item \texttt{audio\_data} ist ein NumPy-Array der Form \texttt{(samples,)} oder \texttt{(samples, channels)}
    \item \texttt{sample\_rate} ist die Abtastfrequenz in Hz
\end{itemize}

\subsubsection{WAV-Export mit Normalisierung}

\begin{lstlisting}
def write_wav(wav_path: str, audio_data, sample_rate: int) -> None:
    """
    Normalizes the audio data to prevent clipping and exports it to a WAV file.
    """
    max_amplitude = np.max(np.abs(audio_data))
    if max_amplitude > 1.0:
        audio_data = audio_data / max_amplitude
    sf.write(wav_path, audio_data, sample_rate)
\end{lstlisting}

\textbf{Algorithmus:}

\begin{algorithm}
    \caption{Normalisierter WAV-Export}
    \begin{algorithmic}
        \STATE $A_{\max} \gets \max_n |y[n]|$
        \IF{$A_{\max} > 1.0$}
            \STATE $y[n] \gets y[n] / A_{\max}$ f\"ur alle $n$
        \ENDIF
        \STATE Schreibe normalisierte Samples in WAV-Datei
    \end{algorithmic}
\end{algorithm}

\subsubsection{Integration in die Pipeline}

\begin{lstlisting}
# main.py Schritt [2/4]: Audio laden
dry_audio, sample_rate = read_wav(input_wav_path)

# Schritt [3/4]: Convolution mit Material-Absorption
conv_material = ir_data.get("metadata", {}).get("wall_material")
wet_audio = convolve(dry_audio, sample_rate, ir_data, material=conv_material)

# Schritt [4/4]: Normalisieren und exportieren
write_wav(output_wav_path, wet_audio, sample_rate)
\end{lstlisting}

Der \texttt{wall\_material}-Wert wird automatisch aus \texttt{ir\_output.json} gelesen, das der Rust-Tracer nach der Simulation schreibt. Damit nutzen Tracer und Convolution stets dasselbe Material.

% ========================================
\subsection{3D-Raumdaten: OBJ-Parser}
% ========================================

\subsubsection{Theoretischer Hintergrund}

\paragraph{OBJ-Format}

Das Wavefront \texttt{.obj}-Format ist ein textbasiertes 3D-Modellformat.
Jede Zeile beginnt mit einem Schl\"usselwort, das den Zeilentyp bestimmt:

\begin{lstlisting}[language={}]
# Vertices (Ecken)
v  1.0  2.0 -1.5    # Vertex 1
v  2.0  2.0 -1.5    # Vertex 2
v  1.0  3.0 -1.5    # Vertex 3

# Faces (Dreiecke, Format: v/vt/vn)
f  1/1/1  2/2/1  3/3/1  # Textur- und Normal-Indizes optional
\end{lstlisting}

Der Parser verarbeitet nur die f\"ur Raytracing relevanten Zeilen:

\begin{table}[H]
    \centering
    \begin{tabular}{|l|l|c|}
        \hline
        \textbf{Zeilentyp} & \textbf{Bedeutung} & \textbf{Verarbeitet?} \\
        \hline
        \texttt{v x y z} & Vertex-Position & \checkmark \\
        \hline
        \texttt{f v1 v2 v3} & Dreiecksfläche (Vertex-Indizes) & \checkmark \\
        \hline
        \texttt{f v/vt/vn} & Face mit Textur- und Normalindex & \checkmark (nur v) \\
        \hline
        \texttt{vn nx ny nz} & Vertex-Normale & Ignoriert \\
        \hline
        \texttt{vt u v} & Textur-Koordinate & Ignoriert \\
        \hline
        \texttt{usemtl name} & Material-Zuweisung (.mtl) & Ignoriert \\
        \hline
        \texttt{\#\ ...} & Kommentar & Ignoriert \\
        \hline
    \end{tabular}
    \caption{OBJ-Zeilentypen und Verarbeitungsstatus im Parser}
\end{table}

\textbf{Indexierung:} OBJ nutzt 1-basierte Indizes, Python 0-basierte.
Der Parser subtrahiert beim Einlesen automatisch 1.

\subsubsection{Dreiecks-Repr\"asentation}

$$\text{Triangle} = \begin{bmatrix}
x_1 & y_1 & z_1 \\
x_2 & y_2 & z_2 \\
x_3 & y_3 & z_3
\end{bmatrix}$$

\subsection{Implementierung: \texttt{obj\_parser.py}}

\subsubsection{OBJ-Loader}

\begin{lstlisting}
def load_obj_to_triangles(file_path: str) -> list:
    vertices = []
    triangles = []
    path = Path(file_path)
    if not path.exists():
        raise FileNotFoundError(f"The 3D file '{file_path}' was not found.")
    with path.open('r') as f:
        for line in f:
            line = line.strip()
            if line.startswith('v '):
                parts = line.split()
                x, y, z = float(parts[1]), float(parts[2]), float(parts[3])
                vertices.append([x, y, z])
            elif line.startswith('f '):
                parts = line.split()
                v1_idx = int(parts[1].split('/')[0]) - 1
                v2_idx = int(parts[2].split('/')[0]) - 1
                v3_idx = int(parts[3].split('/')[0]) - 1
                triangle = [vertices[v1_idx], vertices[v2_idx], vertices[v3_idx]]
                triangles.append(triangle)
    return triangles
\end{lstlisting}

\subsubsection{Material-basierter Export: \texttt{export\_room\_geometry()}}

Die neue Exportfunktion verbindet OBJ-Geometrie mit der Materialdatenbank
(\texttt{materials.json}) und schreibt \texttt{room\_geometry.json} f\"ur den Rust-Tracer.
Jedes Dreieck erh\"alt die mittlere Absorptionskonstante des gew\"ahlten Materials:

\begin{lstlisting}
_MATERIALS_JSON = Path(__file__).resolve().parent.parent / "rust_service" / "materials.json"

def _mean_absorption(material: str) -> float:
    """Mittlere Absorption aus materials.json fuer das gewaehlte Material."""
    data = json.loads(_MATERIALS_JSON.read_text())
    alphas = data["materials"][material]["absorption"]
    return sum(alphas) / len(alphas)

def export_room_geometry(obj_path: str, material: str, out_path: str) -> None:
    triangles_raw = load_obj_to_triangles(obj_path)
    absorption = _mean_absorption(material)
    triangles = [
        {"v0": t[0], "v1": t[1], "v2": t[2], "absorption": absorption}
        for t in triangles_raw
    ]
    out = {"triangles": triangles, "triangle_count": len(triangles)}
    Path(out_path).write_text(json.dumps(out, indent=2))
\end{lstlisting}

\subsubsection{Ausgabeformat: \texttt{room\_geometry.json}}

Der Rust-Tracer liest Dreiecke inkl.\ Material-Absorption aus dieser Datei:

\begin{lstlisting}[language={}]
{
  "triangles": [
    {"v0": [0.0, 0.0, 0.0], "v1": [12.0, 0.0, 0.0],
     "v2": [12.0, 12.0, 0.0], "absorption": 0.033}
  ],
  "triangle_count": 12
}
\end{lstlisting}

Die \texttt{absorption}-Werte entsprechen dem arithmetischen Mittel aller 6 Oktavb\"ander
aus \texttt{materials.json}. Die frequenzabh\"angige Verarbeitung findet ausschlie\ss{}lich
in der Python-Convolution statt.

\subsection{Blender-Workflow: Raum erstellen und exportieren}

F\"ur realistische Raumakustik kann jede beliebige Raumgeometrie aus Blender
exportiert werden. Der empfohlene Workflow:

\begin{enumerate}
    \item \textbf{Raum modellieren}: Geschlossenen Raum in Blender erstellen.
    Alle Fl\"achen m\"ussen nach innen zeigen (\textit{Face Normals} pr\"ufen).
    \item \textbf{Triangulieren}: \textit{Edit Mode} $\to$ \textit{Mesh} $\to$
    \textit{Faces} $\to$ \textit{Triangulate Faces} (Shortcut: \texttt{Ctrl+T}).
    Der OBJ-Parser verarbeitet ausschlie\ss{}lich Dreiecke.
    \item \textbf{Exportieren}: \textit{File} $\to$ \textit{Export} $\to$
    \textit{Wavefront (.obj)}
    \begin{itemize}
        \item Koordinatensystem: Y Forward, Z Up (passend zu \texttt{input\_config.json})
        \item Haken bei \textit{Triangulate Mesh} setzen
        \item Speichern als \texttt{room.obj} in \texttt{python\_service/}
    \end{itemize}
    \item \textbf{Positionen pr\"ufen}: Speaker- und Mic-Positionen in
    \texttt{input\_config.json} m\"ussen innerhalb des Modells liegen.
    Die Mic-Position zuz\"uglich \texttt{mic\_radius} darf keine Wand \"uberschreiten.
    \item \textbf{Pipeline starten}: \texttt{python main.py} erkennt \texttt{room.obj}
    automatisch und generiert \texttt{room\_geometry.json} vor dem Rust-Lauf.
\end{enumerate}

% ========================================
\subsection{Integration und Testing}
% ========================================

\subsubsection{End-to-End-Pipeline}

Die Pipeline verl\"auft nun in 5 Schritten. Schritt~0 ist vollautomatisch:
\texttt{room.obj} wird erkannt, mit dem konfigurierten Material in
\texttt{room\_geometry.json} umgewandelt und dem Rust-Tracer bereitgestellt.

\begin{lstlisting}
# Schritt 0: OBJ -> room_geometry.json (automatisch wenn room.obj vorhanden)
if Path("room.obj").exists():
    export_room_geometry("room.obj", wall_material,
                         "../rust_service/room_geometry.json")

# Schritt 1: Rust Physics Engine
ir_data = run_rust_and_load_json(rust_binary_path, json_path)

# Schritt 2: Audio laden
dry_audio, sample_rate = read_wav(input_wav_path)

# Schritt 3: Convolution mit Material-Absorption
conv_material = ir_data.get("metadata", {}).get("wall_material")
wet_audio = convolve(dry_audio, sample_rate, ir_data, material=conv_material)

# Schritt 4: Normalisieren und exportieren
write_wav(output_wav_path, wet_audio, sample_rate)
\end{lstlisting}

\subsection{Material-Konfiguration}

Das Wandmaterial wird in \texttt{rust\_service/materials.json} festgelegt:

\begin{lstlisting}[language={}]
"materials": {
    "concrete": {
      "absorption": [0.02, 0.03, 0.03, 0.03, 0.04, 0.07],
      "scattering": 0.05
    },
\end{lstlisting}

Vom Rust-Tracer wird \texttt{wall\_material} in \texttt{ir\_output.json} exportiert
und von der Python-Pipeline automatisch f\"ur die Convolution \"ubernommen.
Verf\"ugbare Materialien (aus \texttt{materials.json}):

\begin{table}[H]
    \centering
    \begin{tabular}{|l|l|c|}
        \hline
        \textbf{Material} & \textbf{Beschreibung} & \textbf{Mittl. Absorption} \\
        \hline
        \texttt{concrete} & Beton  & 0.033 \\
        \hline
        \texttt{wood}  & Holz & 0.11 \\
        \hline
        \texttt{glass}    & Glas  & 0.026 \\
        \hline
        \texttt{carpet}   & Teppich  & 0.425 \\
        \hline
    \end{tabular}
    \caption{Verf\"ugbare Materialien in \texttt{materials.json}}
\end{table}

\subsection{Konfiguration: \texttt{input\_config.json}}

Alle Simulationsparameter werden in einer einzigen JSON-Datei geb\"undelt,
die sowohl vom Rust-Tracer als auch von der Python-Pipeline gelesen wird:

\begin{table}[H]
    \centering
    \begin{tabular}{|l|l|p{5.5cm}|}
        \hline
        \textbf{Feld} & \textbf{Beispielwert} & \textbf{Bedeutung} \\
        \hline
        \texttt{speaker\_position} & \texttt{[2.0, 3.0, 1.5]} & 3D-Position der Schallquelle \\
        \hline
        \texttt{mic\_position} & \texttt{[8.0, 7.0, 6.0]} & 3D-Position des Mikrofons \\
        \hline
        \texttt{mic\_radius} & \texttt{3.0} & Einfangradius des Mikrofons in Metern \\
        \hline
        \texttt{rays\_to\_cast} & \texttt{1000} & Anzahl der simulierten Strahlen \\
        \hline
        \texttt{max\_bounces} & \texttt{100} & Max.\ Reflexionen pro Strahl \\
        \hline
        \texttt{min\_pressure} & \texttt{0.001} & Abbruchschwelle (Restdruck) \\
        \hline
        \texttt{wall\_material} & \texttt{"concrete"} & Wandmaterial (aus \texttt{materials.json}) \\
        \hline
    \end{tabular}
    \caption{Felder der \texttt{input\_config.json}}
\end{table}

Das Feld \texttt{wall\_material} wird von der Python-Pipeline beim
Start ausgelesen, um \texttt{room\_geometry.json} mit der passenden Absorption
zu generieren. Der Rust-Tracer schreibt es anschlie\ss{}end in \texttt{ir\_output.json},
von wo die Convolution es f\"ur die frequenzabh\"angige D\"ampfung \"ubernimmt.



% ========================================
\subsection{Visualisierung: \texttt{visualizer.py}}
% ========================================

\subsection{\"Ubersicht}

Der Visualizer rendert die Ergebnisse der Rust-Simulation als interaktive
3D-Szene im Browser (Plotly). Er liest zwei Dateien aus dem \texttt{rust\_service}:

\begin{itemize}
    \item \texttt{visualisation\_data.json}: Strahlpfade, Speaker- und Mic-Position
    \item \texttt{room\_geometry.json}: Raumw\"ande als Dreiecksliste (optional)
\end{itemize}

\subsection{Raumwände-Rendering}

Die Funktion \texttt{\_draw\_room\_walls()} wurde neu implementiert und l\"adt
\texttt{room\_geometry.json}, falls vorhanden, und zeichnet jedes Dreieck als
blaues Drahtgitter:

\begin{lstlisting}
def _draw_room_walls(fig: go.Figure, geometry_path: Path) -> None:
    if not geometry_path.exists():
        return
    data = json.loads(geometry_path.read_text())
    for tri in data["triangles"]:
        v0, v1, v2 = tri["v0"], tri["v1"], tri["v2"]
        xs = [v0[0], v1[0], v2[0], v0[0]]
        ys = [v0[1], v1[1], v2[1], v0[1]]
        zs = [v0[2], v1[2], v2[2], v0[2]]
        fig.add_trace(go.Scatter3d(
            x=xs, y=ys, z=zs, mode='lines',
            line=dict(width=1, color='rgba(100,130,200,0.35)'),
            showlegend=False, hoverinfo='skip'))
\end{lstlisting}

Durch das R\"uckschlie\ss{}en des Dreiecks (\texttt{v0} als vierter Punkt) entsteht
ein geschlossenes Dreieck. Die halbtransparente Farbe \texttt{rgba(100,130,200,0.35)}
h\"alt die Raumw\"ande visuell im Hintergrund, sodass die Strahlpfade im Vordergrund
bleiben.

\subsection{Dargestellte Elemente}

\begin{table}[H]
    \centering
    \begin{tabular}{|l|l|l|}
        \hline
        \textbf{Element} & \textbf{Darstellung} & \textbf{Datenquelle} \\
        \hline
        Strahlpfade & Farbige Linien (Golden-Angle-Farbfolge) & \texttt{visualisation\_data.json} \\
        \hline
        Lautsprecher & Blauer Kreuz-Marker & \texttt{visualisation\_data.json} \\
        \hline
        Mikrofon-Zentrum & Blauer Punkt & \texttt{visualisation\_data.json} \\
        \hline
        Mikrofon-Radius & Rote Kugel (halbtransparent) & \texttt{visualisation\_data.json} \\
        \hline
        Raumw\"ande & Hellblaues Drahtgitter & \texttt{room\_geometry.json} \\
        \hline
    \end{tabular}
    \caption{Visuelle Elemente der 3D-Debugging-Szene}
\end{table}

Die Golden-Angle-Farbfolge $\varphi_i = (i \cdot 137{,}5\,\%) \bmod 360$
verteilt die Strahlfarben gleichm\"a\ss{}ig im HSL-Farbraum, sodass auch bei
vielen Strahlen kein Farbton doppelt vorkommt.

\subsection{Starten des Visualizers}

\begin{lstlisting}[language=bash]
cd python_service
python visualizer.py
\end{lstlisting}

Plotly \"offnet automatisch einen Browser-Tab mit der interaktiven 3D-Szene.
Fehlt \texttt{room\_geometry.json}, werden keine Raumw\"ande gezeichnet;
die Datei wird durch Ausf\"uhren der Pipeline mit einer \texttt{room.obj} erzeugt.

% ========================================
\subsection{Akzeptanzkriterien und Validierung}
% ========================================

\subsubsection{Audio I/O}

\begin{table}[H]
    \centering
    \begin{tabular}{|p{4.5cm}|p{2cm}|p{5cm}|}
        \hline
        \textbf{Kriterium} & \textbf{Erf\"ullt?} & \textbf{Test} \\
        \hline
        Sauberes Einlesen .wav & \checkmark & \texttt{len(audio\_data) > 0} \\
        \hline
        Export ohne Qualit\"atsverlust & \checkmark & H\"ortest, Spektrum-Analyse \\
        \hline
        Normalisierung gegen Clipping & \checkmark & \texttt{max(abs(audio)) <= 1.0} \\
        \hline
    \end{tabular}
\end{table}

\subsubsection{Raumdaten: OBJ-Parser \& Material-Integration}

\begin{table}[H]
    \centering
    \begin{tabular}{|p{4.5cm}|p{2cm}|p{5cm}|}
        \hline
        \textbf{Kriterium} & \textbf{Erf\"ullt?} & \textbf{Test} \\
        \hline
        Parser liest .obj & \checkmark & \texttt{len(triangles) > 0} \\
        \hline
        JSON-Export mit Absorption & \checkmark & \texttt{absorption}-Feld pro Dreieck \\
        \hline
        Rust l\"adt OBJ-Geometrie & \checkmark & \texttt{room\_geometry.json} erkannt \\
        \hline
        Material durchgereicht & \checkmark & \texttt{wall\_material} in \texttt{ir\_output.json} \\
        \hline
        Frequenzabh\"ang. Convolution & \checkmark & \texttt{apply\_material\_absorption()} \\
        \hline
    \end{tabular}
\end{table}

\subsection{Materialvergleich: Akustisches Verhalten}

Die folgende Tabelle zeigt das Verhalten der Pipeline mit dem Test-Raum
(12\,m $\times$ 12\,m $\times$ 10\,m, 100 Strahlen, 100 maximale Bounces):

\begin{table}[H]
    \centering
    \begin{tabular}{|l|c|c|p{5cm}|}
        \hline
        \textbf{Material} & \textbf{Hits} & \textbf{Eff. Bounces} & \textbf{Klangcharakter} \\
        \hline
        \texttt{concrete} & $\sim$2000 & bis 100 & Hell, langer Nachhall (Kathedrale) \\
        \hline
        \texttt{glass}    & $\sim$2200 & bis 100 & Sehr hell, extrem langer Nachhall \\
        \hline
        \texttt{drywall}  & $\sim$1500 & $\sim$43 & Moderat, typischer Wohnraum \\
        \hline
        \texttt{carpet}   & $\sim$400  & $\sim$11 & Dumpf, kurzer Nachhall \\
        \hline
        \texttt{foam}     & $\sim$200  & $\sim$7  & Sehr trocken, fast kein Nachhall \\
        \hline
    \end{tabular}
    \caption[Verhalten verschiedener Materialien]{Verhalten verschiedener Materialien (Richtwerte, 12\,m\,$\times$\,12\,m\,$\times$\,10\,m Testraum)}
\end{table}

Die \textbf{effektiven Bounces} ergeben sich aus der Bedingung
$(1 - \alpha)^n < p_{\min}$ mit $p_{\min} = 0.001$:
$n < \ln(0.001) / \ln(1 - \alpha)$.
Bei Beton ($\alpha = 0.033$) ben\"otigt ein Strahl theoretisch 204 Bounces,
wird aber durch \texttt{max\_bounces = 100} begrenzt.
Bei Teppich ($\alpha = 0.425$) bricht er bereits nach 11 Bounces ab.






